\section{Introduction}
\label{sec:intro}

Reinforcement learning from human feedback (\rlhf) does not optimize what people want; it optimizes a learned reward model (\rmodel) that approximates what people chose in a finite set of comparisons~\citep{christiano2017deep,stiennon2020learning,ouyang2022training}.
When the policy finds features that the \rmodel over-values, measured reward keeps rising while real quality stalls or falls.
This is reward hacking~\citep{amodei2016concrete,skalse2022defining}, and it is a practical problem: a practitioner watching the reward curve cannot tell improvement from exploitation.

Knowing \emph{which} proxy-exploiting behaviors dominate, and \emph{when} they appear, tells practitioners what to monitor and when to stop.
Prior work documents the individual phenomena.
Reward models and \rlhf-trained policies favor longer responses~\citep{singhal2023long}.
Models trained on human feedback tend to agree with users' stated views~\citep{perez2022discovering,sharma2023towards}.
Optimizing far enough against a learned reward lowers true reward~\citep{gao2023scaling,moskovitz2023confronting}.
However, these results come from different models, reward models, and measurement protocols.
What is less common is a single taxonomy applied to the \emph{same} training runs, so that the shares of different failure modes are comparable and each one has an onset time on a shared training axis.

{\bf Which failure modes make up reward hacking, when do they appear, and when does the \rmodel score stop tracking human preference?}
We study this question with a small but complete \rlhf pipeline.
We train an \rmodel on 20,000 human pairwise votes from Chatbot Arena~\citep{chiang2024chatbot}, optimize a \qwensmall policy against it for 160 updates (three seeds plus one strong-KL ablation), and save 11 checkpoints per run.
We score all 10,560 checkpoint responses with the \rmodel, with judge-free surface metrics, and with two blinded LLM judges from different model families that rate quality, compare each response with the untrained policy's response, and assign labels from a six-category taxonomy: verbosity, sycophancy, format gaming, repetition, evasion, and fabrication.
No human raters were available, so the judges are a proxy for human preference throughout.

We fix four hypotheses and their operational definitions before seeing any judge labels:
(H1) verbosity and sycophancy are the primary category in more than 60\% of reward-hacking instances;
(H2) both emerge within the first 20\% of training and then rise monotonically;
(H3) format gaming appears later and saturates;
(H4) human preference diverges from the \rmodel score after 40\% of training.

Most of these fail.
Verbosity plus sycophancy accounts for 49.9\% of reward-hacking instances (95\% CI 45.3--54.3) under the primary judge, and nearly all of that is verbosity: sycophancy is 0.6\%.
The verbosity share is not even an enrichment, because the same judge assigns verbosity as the primary label to 47.1\% of responses from the \emph{untrained} policy.
Verbosity does appear by 20\% of training (+22 tokens per response), but it plateaus after 50--60\% rather than rising monotonically, and sycophancy does not emerge on any of three measures.
Markdown structure is the earliest change, significant at 10\% of training, and it peaks at 50\% and then declines.
The \rmodel score and judge preference are first detectably divergent at the 40\% checkpoint, which roughly matches H4 in timing, but point estimates fall from 20--30\% and judges never clearly prefer the trained policy.
The failure mode that most distinguishes reward-hacking instances is one we did not hypothesize: confident fabrication is the primary label on 34.5\% of instances, against 11.5\% of untrained responses.

In summary, we make the following contributions:
\begin{itemize}[leftmargin=*,itemsep=0pt,topsep=0pt]
    \item We build an end-to-end \rlhf pipeline with dense checkpointing in which one six-category taxonomy, an \rmodel trained on human votes, two LLM judges, and judge-free surface metrics are all applied to the same 10,560 responses.
    \item We test four pre-specified hypotheses about the composition and timing of reward hacking, and report that the composition hypothesis fails (49.9\% against a 60\% threshold) while the timing hypotheses hold only in part.
    \item We show that fabrication, not verbosity or sycophancy, is the category enriched among responses that gain reward but lose the judge, and that label base rates can make a category look dominant when it is not enriched.
    \item We quantify how far LLM-judge taxonomy labels can be trusted: two judges agree on quality (Spearman 0.71) but give estimates of the H1 share that are 30 points apart.
\end{itemize}

These are results about one small \rlhf setup with an LLM-judge proxy for human preference, not about production \rlhf.
\Secref{sec:related} positions our work, \secref{sec:method} describes the pipeline and definitions, \secref{sec:results} reports the results, and \secref{sec:discussion} discusses what they do and do not support.
